% Source: https://artofproblemsolving.com/community/c3h2445360_recurrence_relation_2
% Stable identity: aops-topic-2445360-q01
\textbf{(a)} All terms are positive: $a_1,a_2>0$, and each new term is a quotient of positive numbers. The definition says that
\[a_na_{n-2}-a_{n-1}^2=2\qquad\text{for every }n\ge3.\]
For $n\ge4$, comparing this with the same relation for $n-1$,
\[a_na_{n-2}-a_{n-1}^2=a_{n-1}a_{n-3}-a_{n-2}^2\quad\Longrightarrow\quad a_{n-2}\left(a_n+a_{n-2}\right)=a_{n-1}\left(a_{n-1}+a_{n-3}\right),\]
so
\[\dfrac{a_n+a_{n-2}}{a_{n-1}}=\dfrac{a_{n-1}+a_{n-3}}{a_{n-2}}.\]
Thus the ratio $\dfrac{a_n+a_{n-2}}{a_{n-1}}$ is the same for all $n\ge3$, and it equals its value for $n=3$: $a_3=\dfrac{1+2}{1}=3$ and $\dfrac{a_3+a_1}{a_2}=\dfrac{3+1}{1}=4$. Hence $a_n+a_{n-2}=4a_{n-1}$, i.e. $a_n=4a_{n-1}-a_{n-2}$ for all $n\ge3$. By induction every $a_n$ is an integer. The sequence begins $1,1,3,11,41,153,571,\dots$ $\blacksquare$

\textbf{(b)} The characteristic equation of $a_n=4a_{n-1}-a_{n-2}$ is $r^2=4r-1$, with roots $r=2\pm\sqrt3$. So $a_n=A\left(2+\sqrt3\right)^{n-1}+B\left(2-\sqrt3\right)^{n-1}$ for constants $A,B$, determined by
\[a_1=A+B=1,\qquad a_2=A\left(2+\sqrt3\right)+B\left(2-\sqrt3\right)=1.\]
The second equation gives $2+\sqrt3(A-B)=1$, so $A-B=-\dfrac{1}{\sqrt3}$, and $A=\dfrac{3-\sqrt3}{6}$, $B=\dfrac{3+\sqrt3}{6}$. Therefore
\[\boxed{a_n=\dfrac{\left(3-\sqrt3\right)\left(2+\sqrt3\right)^{n-1}+\left(3+\sqrt3\right)\left(2-\sqrt3\right)^{n-1}}{6}}\qquad(n\ge1).\]
(Check: $n=3$ gives $\dfrac{\left(9+5\sqrt3\right)+\left(9-5\sqrt3\right)}{6}=3$.)
